\chapter{Limiti della corrispondenza}
\label{cap:limiti}

Le grandezze ERA5 della \cref{tab:limiti} non hanno un corrispettivo nell'archivio. Il limite
sta nei file wrfout, che non contengono un campo calcolato dal modello, nella configurazione della
simulazione (\cref{sec:configurazione}) o nella regola di pubblicazione (\cref{sec:pubblicazione}).

\begin{table}[htbp]
  \centering
  \small
  \caption{Grandezze ERA5 assenti dall'archivio e origine del limite.}
  \label{tab:limiti}
  \begin{tabularx}{\textwidth}{@{}XX@{}}
    \toprule
    Parametro & Origine del limite \\
    \midrule
    altezza dello strato limite (\campo{blh}) & file wrfout: \campo{PBLH} calcolato e non scritto \\
    flussi di calore sensibile e latente & file wrfout: \campo{HFX} e \campo{LH} calcolati e non scritti \\
    energia cinetica turbolenta & configurazione: nessuno schema con energia cinetica turbolenta \\
    temperatura massima e minima a \SI{2}{m} (\campo{mx2t}, \campo{mn2t}) & configurazione: \campo{output_diagnostics=0} \\
    onde, coefficiente di resistenza che dipende dallo stato del mare & configurazione: nessun modello d'onda \\
    correnti, altezza della superficie del mare & configurazione: oceano non simulato \\
    spessore e albedo del ghiaccio marino & configurazione: ghiaccio solo come dato di ingresso \\
    precipitazione convettiva (\campo{cp}, \campo{csf}) & configurazione: convezione risolta, campo identicamente nullo \\
    componenti della raffica (\campo{10efg}, \campo{10nfg}) & regola di pubblicazione \\
    indice di area fogliare per classe (\campo{lai_lv}, \campo{lai_hv}) & regola di pubblicazione \\
    angolo e anisotropia dell'orografia sottogriglia (\campo{anor}, \campo{isor}) & regola di pubblicazione \\
    fusione della neve (\campo{smlt}) & regola di pubblicazione \\
    gli 8 livelli di pressione ERA5 sopra \SI{50}{\hPa} (\SIlist{1;2;3;5;7;10;20;30}{\hPa}) & configurazione: top del modello a \SI{50}{\hPa} \\
    \bottomrule
  \end{tabularx}
\end{table}

\section{Campi calcolati e non scritti}
\label{sec:non-scritti}

Lo schema YSU calcola un'altezza dello strato limite (\campo{PBLH}), e con \campo{isfflx=1} i flussi
di calore sensibile e latente (\campo{HFX}, \campo{LH}) sono calcolati a ogni passo temporale. Nessuno
dei tre è stato scritto nei file wrfout, e non esiste un output della simulazione più completo. Lo
stesso vale per la temperatura di pelle del modello \campo{TSK}, che l'archivio ricostruisce dalla
radiazione (\cref{sec:pelle}).

Per l'altezza dello strato limite una diagnosi a posteriori, per esempio con il numero di Richardson
di bulk sui profili, sarebbe nostra e non del modello, e non soddisfa la regola di pubblicazione
(\cref{sec:pubblicazione}). Nemmeno i flussi si ricostruiscono dal bilancio energetico di superficie,
perché nel bilancio restano tre incognite: i due flussi e il flusso di calore nel suolo, anch'esso non
scritto.

\section{Campi che la configurazione esclude}
\label{sec:configurazione-limiti}

Lo schema YSU è non locale e non porta energia cinetica turbolenta, e con \campo{KM_OPT=4} non la
produce nemmeno la diffusione: il campo non esiste a nessuna quota.

Con \campo{output_diagnostics=0} WRF non ha scritto i massimi e i minimi orari della temperatura a
\SI{2}{m}, e un campionamento orario di \campo{2t} non è l'estremo nell'ora; la simulazione porta
massimi orari, con \campo{nwp_diagnostics=1}, ma non per la temperatura.

Nessun modello d'onda è accoppiato all'atmosfera, e mancano quindi l'altezza significativa, il
periodo e la direzione delle onde, e il coefficiente di resistenza che dipende dallo stato del mare.
L'oceano non è simulato (\campo{SF_OCEAN_PHYSICS=0}): non ci sono correnti né altezza della superficie
del mare. Il ghiaccio marino viene solo dai dati di ingresso, come indicatore a due valori per cella,
con albedo costante \num{0.65}: non esistono lo spessore del ghiaccio né un'albedo del ghiaccio che
evolva.

Con la convezione risolta (\campo{CU_PHYSICS=0}) la precipitazione convettiva \campo{RAINC} è
identicamente nulla, e i parametri convettivi di ERA5, \campo{cp} e \campo{csf}, non si pubblicano
(\cref{sec:non-pubblicati}).

\section{Campi che la regola di pubblicazione esclude}
\label{sec:esclusi}

\paragraph{Componenti della raffica.} Per \campo{10efg} e \campo{10nfg} esiste solo il modulo del
massimo orario del vento risolto (\cref{cap:raffica}), e le due componenti sono escluse per
decisione.

\paragraph{Indice di area fogliare.} Il modello porta un solo indice di area fogliare per cella, un
valore di tabella per categoria dominante (\campo{rdlai2d=.false.}), mentre ERA5 ne vuole uno per la
vegetazione bassa e uno per quella alta. Qualunque divisione del valore unico fra le due sarebbe
un'ipotesi nostra. Il campo \campo{LAI} del modello è nei file wrfout e non è pubblicato.

\paragraph{Angolo e anisotropia dell'orografia sottogriglia.} Il file statico porta i parametri di
asimmetria (\campo{OA1-4}) e di lunghezza efficace (\campo{OL1-4}) dello schema di Kim--Arakawa.
\campo{anor} e \campo{isor} di ECMWF vengono dalla struttura della covarianza del gradiente del
terreno, e le due grandezze sono diverse.

\paragraph{Fusione della neve.} Il campo \campo{ACSNOM} esiste ma non è un accumulatore monotono. In
15 dei 16 giorni esaminati alcuni punti calano fra le 00 e le 12 UTC, fino a \SI{8.2}{mm}, e il 1º
luglio 2024 il calo rispetto alle 00 UTC arriva a \SI{7.26}{mm}. Costruito come gli altri
campi accumulati, come differenza fra due ore successive della stessa run (\cref{sec:tempo}),
\campo{smlt} sarebbe negativo in quei punti.

\section{Campi presenti e non pubblicati}
\label{sec:non-pubblicati}

Sedici campi dei file wrfout sono identicamente nulli in tutti gli istanti del campione di verifica:
\campo{ACHFX}, \campo{ACLHF}, \campo{ACGRDFLX}, \campo{NOAHRES}, \campo{SSTSK}, \campo{SST_INPUT},
\campo{SWNORM}, \campo{ACSNOW}, \campo{HAILNC}, \campo{RAINC}, \campo{RAINSH}, \campo{PREC_ACC_C} e i
quattro campi della radiazione termica discendente al top dell'atmosfera, \campo{LWDNT},
\campo{LWDNTC}, \campo{ACLWDNT} e \campo{ACLWDNTC}. \campo{ACHFX}, \campo{ACLHF}, \campo{ACGRDFLX} e
\campo{NOAHRES} sono nulli perché sono accumulatori dello schema di superficie Noah, e questa
simulazione usa RUC. Un campo identicamente nullo non si pubblica, qualunque sia la ragione.

Per altri campi non nulli nessun parametro ECMWF descrive la stessa grandezza: la frazione congelata
della precipitazione (\campo{SR}), la densità della neve che cade (\campo{RHOSNF}), i massimi e minimi
annuali della frazione di vegetazione verde (\campo{SHDMAX}, \campo{SHDMIN}) e la temperatura del suolo
profondo usata come condizione al contorno inferiore (\campo{TMN}). Il deflusso accumulato
\campo{ACRUNOFF} coincide con il deflusso superficiale \campo{SFROFF}, ed è già in \campo{sro}.

Cinque campi non nulli restano fuori anche se ECMWF definisce un parametro per la stessa grandezza,
perché ERA5 non lo pubblica e l'archivio segue lo schema di ERA5 (\cref{tab:non-pubblicati}).

\begin{table}[htbp]
  \centering
  \small
  \caption{Campi non nulli dei file wrfout con un parametro ECMWF che ERA5 non pubblica.}
  \label{tab:non-pubblicati}
  \begin{tabularx}{\textwidth}{@{}lX@{}}
    \toprule
    Campo WRF & Parametro ECMWF \\
    \midrule
    \campo{SNOALB}     & 260161 \campo{mxsalb}, albedo massima della neve \\
    \campo{QGRAUP}     & 260028 \campo{grle} e 260001 \campo{tcolg}, graupel e sua colonna (\cref{sec:colonne}) \\
    \campo{REFL_10CM}  & 260134 \campo{bref}, riflettività radar \\
    \campo{LPI}        & famiglia 228050 \campo{litoti}; \campo{LPI} è però un indice di potenziale di fulminazione, non una densità di fulmini \\
    \campo{SNOWFALLAC} & 3066 \campo{sde}, perché è uno spessore geometrico della neve e non un'acqua equivalente \\
    \bottomrule
  \end{tabularx}
\end{table}
